% !TEX program = xelatex
% 공통수학2 · Ⅱ. 집합과 명제 · 1. 집합 · 04 여집합과 차집합 · 1/2차시
% 학생용 활동지 (답안 없음)
\documentclass[11pt,a4paper]{article}
\usepackage{kotex}
\usepackage{iftex}
\ifPDFTeX\else
  \setmainhangulfont{Noto Serif CJK KR}
  \setsanshangulfont{Noto Sans CJK KR}
\fi
\usepackage[margin=2.1cm]{geometry}
\usepackage{parskip}
\usepackage{enumitem}
\usepackage{array}
\usepackage{tabularx}
\usepackage{xcolor}
\usepackage{amssymb}
\usepackage{tikz}
\usepackage[most]{tcolorbox}
\usepackage{fancyhdr}
\usepackage{titlesec}

\definecolor{goldc}{HTML}{B07C22}
\definecolor{goldstrong}{HTML}{8A5F16}
\definecolor{indigoc}{HTML}{2B3B57}
\definecolor{indigosoft}{HTML}{6E82A6}
\definecolor{linec}{HTML}{D6DACB}
\definecolor{surface2}{HTML}{E4E8DC}
\definecolor{writeline}{HTML}{C7CCBA}
\definecolor{inksoft}{HTML}{52604F}

\titleformat{\section}{\normalfont\sffamily\bfseries\large\color{indigoc}}{\thesection}{0.6em}{}
\titlespacing{\section}{0pt}{14pt}{6pt}
\newtcolorbox{goalbox}{enhanced, breakable, colback=white, colframe=goldc, boxrule=0.5pt, leftrule=3pt, arc=2pt, left=10pt, right=10pt, top=8pt, bottom=8pt}
\newtcolorbox{sheetbox}{enhanced, breakable, colback=white, colframe=linec, boxrule=0.6pt, arc=3pt, left=11pt, right=11pt, top=9pt, bottom=9pt}
\newcommand{\writespace}[1]{%
  \par\nobreak\vspace{3pt}
  \foreach \n in {1,...,#1} {%
    \noindent\textcolor{writeline}{\rule{\linewidth}{0.4pt}}\par\vspace{0.62cm}%
  }
}
\newcommand{\fillblank}[1]{\makebox[#1]{\hrulefill}}
\pagestyle{fancy}\fancyhf{}\renewcommand{\headrulewidth}{0pt}
\fancyfoot[L]{\footnotesize\color{inksoft} 공통수학2 · 2026학년도 2학기 · 지평선고등학교}
\fancyfoot[R]{\footnotesize\color{inksoft} 다음 시간 --- 04 여집합과 차집합(2)}

\begin{document}

\noindent
\begin{minipage}[b]{0.62\linewidth}
  {\sffamily\footnotesize\color{indigosoft} 공통수학2 · Ⅱ. 집합과 명제 · 1. 집합}\\[2pt]
  {\sffamily\bfseries\Large 04 여집합과 차집합 \textperiodcentered\ 23차시}
\end{minipage}\hfill
\begin{minipage}[b]{0.36\linewidth}
  \raggedleft\small\color{inksoft}
  반\ \fillblank{1.6cm} \quad 번호\ \fillblank{1.6cm}\\[4pt]
  이름\ \fillblank{3.6cm}
\end{minipage}

\vspace{4pt}
\noindent{\color{goldc}\rule{\linewidth}{2.2pt}}
\vspace{10pt}

\begin{goalbox}
\textbf{오늘의 목표}\\
전체집합과 여집합, 차집합의 뜻을 알고 구할 수 있다.
\vspace{4pt}
\small\color{inksoft}
$\square$ 자신 있음 \quad $\square$ 조금 헷갈림 \quad $\square$ 처음 봄 \hfill (수업 전 나의 상태에 표시)
\end{goalbox}

\section{생각 열기 --- 태양계 행성}

\begin{sheetbox}
태양계 행성 전체의 집합 $U=\{$수성, 금성, 지구, 화성, 목성, 토성, 천왕성, 해왕성$\}$. 이 중 지구보다 태양에 더 가까이 있는 행성의 집합을 $A$라 하자.

\vspace{4pt}
집합 $U$의 원소 중에서 집합 $A$에 속하지 않는 원소를 모두 말해 보자.
\writespace{2}
\end{sheetbox}

\section{개념 정리 --- 전체집합·여집합·차집합}

\begin{sheetbox}
어떤 집합의 부분집합을 생각할 때, 처음의 집합을 \textbf{전체집합}이라 하며 기호 $U$로 나타낸다.

\vspace{6pt}
전체집합 $U$의 부분집합 $A$에 대하여, $U$의 원소 중 $A$에 속하지 않는 모든 원소로 이루어진 집합을 $A$의 \textbf{여집합}이라 하고 $A^C = \{x\mid x\in U$ 그리고 $x\ \fillblank{1.4cm}\ A\}$

\vspace{6pt}
두 집합 $A$, $B$에 대하여 $A$에 속하지만 $B$에는 속하지 않는 모든 원소로 이루어진 집합을 $A$에 대한 $B$의 \textbf{차집합}이라 하고 $A-B = \{x\mid x\in A$ 그리고 $x\ \fillblank{1.4cm}\ B\}$

\vspace{8pt}
{\footnotesize\color{inksoft} 주의: $A-B$는 $A$에서 $B$를 빼는 것이 아니라, $A$와 $B$의 \fillblank{4cm}을 $A$에서 제외하는 것이다.}
\end{sheetbox}

\section{연습 문제}

\begin{sheetbox}
\textbf{1.} 전체집합 $U=\{x\mid x$는 10 이하의 자연수$\}$에 대하여 다음 집합의 여집합을 구하시오.\\
\hspace*{1.6em}⑴ $A=\{1,2,3,6\}$ \qquad ⑵ $B=\{x\mid x$는 9의 약수$\}$
\writespace{3}

\vspace{6pt}
\textbf{2.} 다음 두 집합 $A$, $B$에 대하여 $A-B$와 $B-A$를 구하시오.\\
\hspace*{1.6em}⑴ $A=\{2,3,5,7,9\}$, $B=\{3,9\}$\\
\hspace*{1.6em}⑵ $A=\{x\mid x$는 18의 약수$\}$, $B=\{x\mid x$는 12의 약수$\}$
\writespace{4}

\vspace{6pt}
\textbf{3.} 전체집합 $U$의 부분집합 $A$에 대하여 다음이 성립함을 벤 다이어그램을 그려 확인하시오.\\
\hspace*{1.6em}⑴ $(A^C)^C=A$ \qquad ⑵ $A\cup A^C=U$ \qquad ⑶ $A\cap A^C=\varnothing$
\writespace{4}
\end{sheetbox}

\section{사고의 가시화 --- See · Think · Wonder}

\noindent{\footnotesize\color{inksoft} 오늘 배운 `여집합과 차집합'을 떠올리며 아래 세 칸을 채워 보자.}\\[6pt]
\begin{tabularx}{\linewidth}{@{}XXX@{}}
\textbf{\footnotesize\color{indigoc} See --- 무엇을 보았나?} &
\textbf{\footnotesize\color{indigoc} Think --- 무엇을 생각했나?} &
\textbf{\footnotesize\color{indigoc} Wonder --- 무엇이 궁금한가?} \\[4pt]
\end{tabularx}
\noindent
\begin{minipage}[t]{0.32\linewidth}\writespace{3}\end{minipage}\hfill
\begin{minipage}[t]{0.32\linewidth}\writespace{3}\end{minipage}\hfill
\begin{minipage}[t]{0.32\linewidth}\writespace{3}\end{minipage}

\section{마무리 성찰}

\begin{sheetbox}
\begin{minipage}[t]{0.48\linewidth}
{\footnotesize\color{inksoft} 오늘 배운 것을 한 문장으로}
\writespace{2}
\end{minipage}\hfill
\begin{minipage}[t]{0.48\linewidth}
{\footnotesize\color{inksoft} 감사 일기 / 유념 한 줄}
\writespace{2}
\end{minipage}
\end{sheetbox}

\end{document}
